\section{Exact algebraic decision procedure}
\label{sec4}\label{sec:decision}

The compressed criterion of Section~\ref{sec:prelim} is decided by a greatest-common-divisor test on
the maximal minors. The closedness of the bad locus, the generic-point criterion and the finite-cover
principle that accompany it are results of \cite{companion1}, recalled below as black boxes. What they
do not supply, and what this section develops, is the passage from a \emph{generic} point to an
\emph{entire chart}, so that one computation per chart gives a verdict valid at every point of it. In
the hyperplane layer that passage is supplied by a single invariant, $\rank H\equiv\omega$, which keeps
a common factor from degenerating to a constant (Lemma~\ref{sec4:lem:fullrank}) --- the counterpart of
its inherited use, which keeps common roots from escaping to infinity. No complexity assertion is made
here; that is the content of Section~\ref{sec5}.

\subsection{Standing notation}
\label{sec4:notation}

We follow Sections~\ref{sec:prelim} and~\ref{sec3}. Let $\mathcal L=\binom{L_0}{R}$ have $q$ rows and
$V=\row\mathcal L$, let $V_0=\row L_0$, $N$ be a basis of $K_0$, and
\[
\begin{gathered}
H=LN,\qquad H_A=LAN,\\
P(\lambda)=\lambda H-H_A,\qquad \omega=\rank H ,
\end{gathered}
\]
with $\omega=k-1$ in the hyperplane layer, $\rho=\rank[C;CA]$, $s=q-k+1$, and
$\widetilde r=r(v)-\delta$ the rank of the quotient bundle. A parameter point $X$ of a chart gives
$L(X)$, $H(X)=L(X)N$, $H_A(X)=L(X)AN$ and $P(X,\lambda)=\lambda H(X)-H_A(X)$, a $q\times k$ pencil of
rank $\omega$ at every $X$. The index $I$ ranges over pairs of $\omega$ rows and $\omega$ columns, and
$m_I(X,\lambda)$ is the corresponding maximal minor.

\subsection{Inherited tools}
\label{sec4:inherited}

\begin{thm}[Inherited tools; \cite{companion1}]
\label{sec4:thm:closed}
\label{sec4:thm:generic}
\label{sec4:prop:cover}
Let the layer be covered by finitely many irreducible parameter pieces with $\rank H\equiv\omega$ and
function field $K=\overline{\mathbb Q(\mathcal S)}$. Then:
\begin{enumerate}[label=(\alph*),leftmargin=*]
\item \label{it:closed} for every chart $\mathcal X$ the bad locus $\mathcal B\subseteq\mathcal X$ is
Zariski closed;
\item \label{it:generic} a piece $\mathcal S$ contains a good point if and only if
$\deg_\lambda\gcd_I m_I=0$ in $K[\lambda]$;
\item \label{it:cover} consequently $\mathcal X$ contains a good point if and only if some piece of the
cover has a gcd of $\lambda$-degree zero.
\end{enumerate}
\end{thm}

Closedness in \cite{companion1} is obtained as follows. Homogenize $\lambda=u/v$ to
$[u{:}v]\in\mathbb P^1$ and let
$\mathcal I=\{(X,[u{:}v]):\ m_I(X,u,v)=0\ \text{for all}\ I\}\subseteq\mathcal X\times\mathbb P^1$ be
the incidence set of common roots, where $m_I(X,u,v)$ is the homogenization of $m_I$. Since every
$m_I$ is polynomial in $(X,u,v)$ and equations are closed, $\mathcal I$ is closed; its image under the
projection $\mathcal X\times\mathbb P^1\to\mathcal X$ is closed because $\mathbb P^1$ is proper
\cite{bpr06}; and a common root at infinity, $[u{:}v]=[1{:}0]$, would force every leading coefficient,
i.e.\ every $\omega\times\omega$ minor of $H(X)$, to vanish at $X$, which $\rank H=\omega$ excludes.
Hence the image of $\mathcal I$ is exactly the locus where a finite common root exists, which is the
bad locus $\mathcal B$. If $\mathcal X$ is irreducible, $\mathcal B$ is closed and distinct from
$\mathcal X$ unless the gcd over the function field has positive degree; a nonzero $\gcd$ of positive
degree gives a common factor, and a Bézout identity over the function field gives the converse
\cite{lan02}, which is (b). The finite-cover statement (c) is the specialization to a finite union.
Two consequences are used below: no decomposition into irreducible components is needed, because the
charts of Section~\ref{sec4:charts} are affine spaces covering the layer by
Theorem~\ref{sec3:thm:param}; and the criterion speaks of the \emph{generic} point only, which is the
gap that Lemma~\ref{sec4:lem:fullrank} closes.

\subsection{Charts and their size}
\label{sec4:charts}

\begin{defn}[Good points, and affine charts of the fibre]
\label{sec4:defn:goodbad}
\label{sec4:defn:charts}
A parameter point $X$ is \emph{good} if $\rank P(X,\lambda)=\omega$ for every $\lambda\in\mathbb C$,
and \emph{bad} otherwise; by the criterion of Section~\ref{sec:prelim} it is bad exactly when
$\deg_\lambda\gcd_I m_I(X,\cdot)>0$. The charts are constructed as follows. Fix an admissible direction
$[v]\in\mathbb P(E_\mu^0)$ and a basis $e_1,\dots,e_{\widetilde r}$ of
$\mathcal Q_v/\overline V_{0,v}$. For $I\subseteq\{1,\dots,\widetilde r\}$ with $|I|=m=s-\delta$ put
$J=\{1,\dots,\widetilde r\}\setminus I$ and let
$\mathcal X_I\cong\mathbb C^{m\times(\widetilde r-m)}$ be the affine chart with coordinates
$X=(X_{ij})_{i\in I,\,j\in J}$. The row space attached to $X$ is spanned by the base rows --- a basis
of $V_0+D_v$, of which there are $(k-1)+\delta$ --- together with
\[
\row_i=e_i+\sum_{j\in J}X_{ij}\,e_j,\qquad i\in I .
\]
These are $(k-1)+\delta+m=(k-1)+s=q$ rows and they form $L(X)$. The chart count is
$\binom{\widetilde r}{m}$, the parameter count is
$d=m(\widetilde r-m)=(s-\delta)(r(v)-s)\le p^2$, and for $m=0$ the chart is a single point.
\end{defn}

The charts above are \emph{affine} and \emph{explicitly} parameterized: the rows are affine in $X$, so
the minors are polynomials a computer can form, and each chart is irreducible by construction, so
Theorem~\ref{sec4:thm:generic} applies without any decomposition. What they do not yet provide is a
cover of the \emph{total} parameter space: when $\dim E_\mu^0>1$ the direction is itself a parameter and
the rows of $L$ depend on it through $D_v$ and through a basis of the quotient. The next definition
supplies the missing charts, so that the algorithm of Section~\ref{sec4:algorithm} enumerates a finite
list rather than the points of a continuum.

\begin{defn}[Charts of the total parameter space]
\label{sec4:defn:totcharts}
Fix a real eigenvalue $\mu$ with $K=\mathbb Q(\mu)$ and the real embedding selecting $\mu$, let
$d_\mu^0=\dim E_\mu^0$, $h_\mu^0=\dim(E_\mu^0\cap\ker A_{21})$, and fix a $K$-basis
$e_1,\dots,e_{d_\mu^0}$ of $E_\mu^0$. A chart of the total parameter space is a product of three
factors, each a principal open subset of an affine space over $K$.
\begin{enumerate}[label=(\roman*),leftmargin=*]
\item \emph{Direction chart.} If $d_\mu^0=1$ the factor is a point. Otherwise, for each
$j\in\{1,\dots,d_\mu^0\}$ let $v_j(x)=e_j+\sum_{i\ne j}x_i e_i$ with coordinates $x\in K^{d_\mu^0-1}$;
these $d_\mu^0$ charts cover $\mathbb P(E_\mu^0)$. The eigen stratum
$\mathbb P(E_\mu^0\cap\ker A_{21})$ is a projective subspace of dimension $h_\mu^0-1$; it is covered by
the $h_\mu^0$ charts obtained in the same way from a $K$-basis of $E_\mu^0\cap\ker A_{21}$. The non-eigen
stratum is covered by the intersections of the $v_j$-charts with the opens
$\{(A_{21}v_j(x))_q\ne0\}$, one for each row index $q$ of $A_{21}$; the factors
$\{(A_{21}v_j(x))_q\}$ are the pivot conditions of the trivialization.
\item \emph{Trivialization chart.} On each direction chart the two functionals $\varphi_v,\varphi_{Av}$
have entries that are affine in $x$, so $D_v=K_0\cap U_v^\perp$, $U_v^\perp$ and the quotient
$\mathcal Q_v/\overline V_{0,v}$ are obtained by Gaussian elimination over $K(x)$; the pivots are
specified minors of the $2\times k$ matrix $[\varphi_v;\varphi_{Av}]$ and of the analogous
$2\times n$ matrix, and each choice of pivots gives a basis whose entries are rational functions of $x$
with denominators equal to those minors. Refining the direction chart by those nonvanishing conditions
gives a locally closed affine chart on which the basis is polynomial after clearing denominators.
\item \emph{Fibre chart and rank chart.} On the trivialized chart, add a Grassmann chart of
Definition~\ref{sec4:defn:charts} for the fibre, and refine by the condition that a specified
$\omega\times\omega$ minor $\Delta$ of $H$ is nonzero; on that open set $\rank H=\omega$ identically.
\end{enumerate}
\end{defn}

\begin{lem}[The total charts cover, and are finitely many]\label{sec4:lem:cover}
Every admissible point of the layer lies in at least one chart of Definition~\ref{sec4:defn:totcharts};
the number of charts is finite, and on each of them the entries of $L$, $H$ and $H_A$ are rational
functions of the joint coordinates whose denominators are the pivot and rank minors, while
$\rank H\equiv\omega$ holds identically.
\end{lem}

\begin{proof}
An admissible point determines a real eigenvalue $\mu$, a direction $[v]\in\mathbb P(E_\mu^0)$, a point
of the fibre over $[v]$ and a positive value of $\rank H$ equal to $\omega$; the direction lies in one
of the direction charts, the elimination pivots are those of its own trivialization, some
$\omega$-minor $\Delta$ of $H$ is nonzero there, and the fibre point lies in one Grassmann chart. Hence
the point lies in the corresponding product chart. Finiteness follows because each of the three factors
is indexed by a finite set: $d_\mu^0$ (or $h_\mu^0$) projective charts, finitely many pivot choices,
finitely many $\omega$-minors and $\binom{\widetilde r}{m}$ Grassmann charts, and there are at most $k$
real eigenvalues. The rationality statement is the content of (ii) and (iii): the elimination is
performed over $K(x)$, so the basis entries are ratios of minors, and clearing denominators makes them
polynomial on the refined open set; $\rank H=\omega$ is imposed by $\Delta\ne0$.
\end{proof}

\begin{lem}[Size of a chart]
\label{sec4:lem:minorcount}
On every chart of Definition~\ref{sec4:defn:charts} the entries of $L(X)$, hence of $H(X)$ and
$H_A(X)$, are affine in $X$; every maximal minor satisfies $\deg_X m_I\le\omega$ and
$\deg_\lambda m_I\le\omega$; and the number of maximal minors is
$\binom{q}{\omega}\binom{k}{\omega}=\binom{q}{k-1}k=k\binom{q}{s}$.
\end{lem}

\begin{proof}
$H(X)=L(X)N$ is affine in $X$ because $N$ is constant, and $m_I$ is an $\omega\times\omega$
determinant, hence multilinear, so its degree is at most $\omega$ in each of $X$ and $\lambda$. The
count is the number of ways of choosing $\omega$ of the $q$ rows and $\omega$ of the $k$ columns,
with $\binom{k}{k-1}=k$ and $\binom{q}{k-1}=\binom{q}{s}$.
\end{proof}

\subsection{The full-rank leading-coefficient lemma}
\label{sec4:fullrank}

\emph{Why a generic test is not yet a decision.} Theorem~\ref{sec4:thm:generic} tests the generic point
of an irreducible piece, which answers a question about almost every point, whereas a decision
procedure needs an answer at every point: a chart whose generic point is bad may in principle contain
good points. The gap is real for general families; what closes it here is the following, the main new
result of this section.

The coefficient of $\lambda^\omega$ in $m_I(X,\lambda)$ is the $\omega\times\omega$ minor of $H(X)$
with the same index set. Since $\rank H(X)=\omega$ at every point of a chart, at least one of these
minors is nonzero, that is,
\begin{equation}
\label{sec4:eq:fullrank}
\text{for every }X\in\mathcal X_I\ \text{there is }I\text{ with }
\deg_\lambda m_I(X,\cdot)=\omega .
\end{equation}

\begin{lem}[Full-rank leading coefficient; dichotomy on a chart]
\label{sec4:lem:fullrank}
\label{sec4:cor:dichotomy}
Let $K$ be a field, $X=(X_1,\dots,X_d)$, and let $m_1,\dots,m_t\in K[X,\lambda]$ satisfy
$\deg_\lambda m_i\le\omega$ and \eqref{sec4:eq:fullrank}. If
$\deg_\lambda\gcd(m_1,\dots,m_t)>0$ in $K(X)[\lambda]$, then for \emph{every} $X_0$ the specialized
polynomials $m_i(X_0,\cdot)$ have a nonconstant common factor. In particular, on a chart of the layer
exactly one of the following holds: either $\deg_\lambda\gcd_I m_I=0$, in which case the chart contains
a good point and its good points form a nonempty Zariski open subset; or
$\deg_\lambda\gcd_I m_I>0$, in which case \emph{every} point of the chart is bad.
\end{lem}

\begin{proof}
Let $g$ be the gcd in $K(X)[\lambda]$, of positive degree. By Gauss's lemma \cite{lan02} we may take
$g$ primitive in $K[X][\lambda]$ and of positive degree, with $g\mid m_i$ in $K[X][\lambda]$ for every
$i$; write $m_i=g\,h_i$ and $\omega_i=\deg_\lambda m_i$, $\gamma=\deg_\lambda g\ge1$. Fix $X_0$. If
$g(X_0,\cdot)\equiv0$ then $m_i(X_0,\cdot)=g(X_0,\cdot)h_i(X_0,\cdot)\equiv0$ for every $i$,
contradicting \eqref{sec4:eq:fullrank}. If $g(X_0,\cdot)$ is a nonzero constant then
$\deg_\lambda m_i(X_0,\cdot)\le\deg_\lambda h_i\le\omega_i-\gamma\le\omega-\gamma<\omega$ for every
$i$, again contradicting \eqref{sec4:eq:fullrank}. Hence $g(X_0,\cdot)$ is nonzero and nonconstant, and
it is a common factor of all $m_i(X_0,\cdot)$. The dichotomy follows by applying this to the charts and
Theorem~\ref{sec4:thm:generic} to the degree-zero alternative.
\end{proof}

Without \eqref{sec4:eq:fullrank} the implication is false: $m(X,\lambda)=X\lambda+1$ has gcd $m$, of
degree $1>0$ over $K(X)[\lambda]$, yet at $X_0=0$ it specializes to the constant $1$, and
\eqref{sec4:eq:fullrank} fails there. The lemma thus rests on the invariant $\rank H\equiv\omega$,
which is what the layer supplies.

\subsection{Real designs and certificates}
\label{sec4:real}\label{sec4:certificates}

\begin{prop}[Real designs and certificates]
\label{sec4:prop:realpoints}
\label{sec4:prop:sieve}
\label{sec4:prop:dimineq}
\label{sec4:prop:resultant}
\label{sec4:prop:cert}
The following hold, in decreasing logical strength.
\begin{enumerate}[label=(\alph*),leftmargin=*]
\item \label{it:real} Every chart $\mathcal X_I\cong\mathbb C^d$ carries the standard real structure
$\mathbb R^d$, and so do the affine charts of $\mathbb P(E_\mu^0)$ because $\mu$ is real and $E_\mu^0$
is defined by real linear equations. The real points of every chart are Zariski dense in it; if the
good set of a chart is a nonempty Zariski open subset then it contains a real point, hence a real
admissible design; and if $\deg_\lambda\gcd>0$ on a chart then every real point of it is bad.
\item \label{it:sieve} If $E_\mu^0=\{0\}$ for every real eigenvalue $\mu$ of
$M=\pi_{K_0}A|_{K_0}$, then the layer is infeasible.
\item \label{it:dimineq} If the fibre over some $[v]\in\mathbb P(E_\mu^0)$ is nonempty then
$\delta\le s\le r(v)$; if $\delta>s$, or $s>\rho$ on the eigen branch and $s>\rho-1$ on the non-eigen
branch, the layer is infeasible.
\item \label{it:resultant} If $\operatorname{Res}_\lambda(m_I,m_J)\not\equiv0$ in $K[X]$ for some pair
$I,J$, then $\deg_\lambda\gcd_I m_I=0$ and the chart, hence the layer, is feasible. The converse fails:
$f_1=\lambda(\lambda-1)$, $f_2=\lambda(\lambda-2)$ and $f_3=(\lambda-1)(\lambda-2)$ share a factor
pairwise, respectively $\lambda$, $\lambda-1$ and $\lambda-2$, yet $\gcd(f_1,f_2,f_3)=1$; a nonzero
resultant is a usable positive certificate, whereas failure to find one is not a negative certificate.
\item \label{it:cert} A complete decision certificate consists of (i) the factorization of $\chi_M$ and
a basis of $E_\mu^0$ over each irreducible factor, or the statement $E_\mu^0=0$; (ii) for every
admissible direction and chart, the chart and parameter counts; (iii) for every chart either a
\emph{good} alternative --- a Bézout identity $\sum_I a_I m_I=1$ with $a_I\in K[X][\lambda]$, checked
by one multiplication --- or a \emph{bad} alternative --- a nonconstant $g\in K[X][\lambda]$ with
division identities $g\,h_I=m_I$ for every $I$. If some chart reports ``good'' the layer is feasible;
if every chart reports ``bad'' it is infeasible.
\end{enumerate}
\end{prop}

\begin{proof}
(a) The zero set of a nonzero polynomial has empty interior in $\mathbb R^d$, so its complement is
dense \cite{bpr06}; the good set is a non-vanishing set of a nonzero polynomial, so it meets
$\mathbb R^d$, and the corresponding $L(X)$ is a real matrix with a real row space.
(b) Every admissible direction lies in some $E_\mu^0$.
(c) The fibre is $\Gr(m,\widetilde r)$ with $m=s-\delta$, nonempty exactly when
$0\le m\le\widetilde r$.
(d) A nonzero resultant makes $m_I$ and $m_J$ coprime in $K(X)[\lambda]$, so the gcd is $1$; the
failure of the converse is the example displayed.
(e) The good alternative certifies $\gcd=1$ without any gcd computation; for the bad alternative,
Lemma~\ref{sec4:lem:fullrank} shows that $g$ is nonconstant and divides every minor at every point, so
the whole chart is bad, and the cover is finite.
\end{proof}

\subsection{The decision algorithm}
\label{sec4:algorithm}

\begin{algo}[Exact decision for the entire layer]
\label{sec4:algo:decide}
\begin{enumerate}
\item \emph{Pre-check.} If $\rank L_0<r_0$, output ``infeasible'' and stop: $\mathcal L$ must have
full row rank, so $\rank L_0=r_0$ is necessary.
\item Compute $K_0$, $N$ and $M=\pi_{K_0}A|_{K_0}$; factor $\chi_M$ and keep its real irreducible
factors.
\item For every real irreducible factor, enumerate the charts of Definition~\ref{sec4:defn:totcharts}:
first the direction charts of \ref{sec4:defn:totcharts}(i), taking the eigen charts when
$h_\mu^0\ge1$ and the non-eigen charts when $h_\mu^0<d_\mu^0$, and compute $\delta,s,m,\widetilde r$ on
each; if $\delta>s$ or $s>r(v)$, skip that chart by Proposition~\ref{sec4:prop:dimineq}. This
enumeration is finite by Lemma~\ref{sec4:lem:cover} and does \emph{not} sample individual directions.
\item For every chart produced in step 3: trivialize the quotient as in
Definition~\ref{sec4:defn:totcharts}(ii), attach a fibre chart \ref{sec4:defn:totcharts}(iii), build
$L(X)$ over $K(x)$ with $x$ the joint coordinates, clear the denominators of the pivot and rank minors,
form the $k\binom{q}{s}$ maximal minors of $P(X,\lambda)$, and compute their gcd over $K(x)[\lambda]$
by the subresultant algorithm \cite{vzgg03}.
\item If some chart has $\deg_\lambda\gcd=0$, output ``feasible'' and stop.
\item If every chart has $\deg_\lambda\gcd>0$, output ``infeasible''.
\item If the computation cannot be completed, output ``undecided'', never ``infeasible''.
\end{enumerate}
\end{algo}

\begin{prop}[Correctness]
\label{sec4:prop:correct}
Whenever steps 1--6 are completed, the verdict of Algorithm~\ref{sec4:algo:decide} is correct.
\end{prop}

\begin{proof}
Step 1 is the problem definition and step 3 is Proposition~\ref{sec4:prop:dimineq}. Step 3 is finite by
Lemma~\ref{sec4:lem:cover}, and its charts cover the whole parameter space, so no admissible design is
lost. In step 5 a chart with $\deg=0$ contains a good point by Theorem~\ref{sec4:thm:generic}; in step 6
Lemma~\ref{sec4:cor:dichotomy} makes every chart entirely bad, and the cover is finite, so no good
point exists. The ``undecided'' output of step 7 is not a verdict. Two implementation points belong to
the proof as well: step 2 enumerates irreducible factors rather than roots, since eigenspaces of
conjugate roots are Galois conjugate and conjugation preserves the $\lambda$-degree of the gcd; and
all linear algebra in step 4 must be done over the \emph{same} field $K=\mathbb Q(\mu)$, because when
$\mu\notin\mathbb Q$ the bases of $E_\mu^0$, $D_v$, $U_v^\perp$ and
$\mathcal Q_v/\overline V_{0,v}$ have no rational coordinates, and one coefficient field is what makes
the chain ``chart rows, minors, function-field gcd'' correct.
\end{proof}

Three statements of different strength are kept apart. \textbf{(i)} This section gives an exact
decision for the entire layer and makes no complexity assertion. \textbf{(ii)} Polynomial time on the
geometrically simple subclass belongs to Section~\ref{sec5}, and an input outside it can only be
answered ``undecided''. \textbf{(iii)} The counterexample of Section~\ref{sec:counterexample} refutes
the reduction ``one representative direction per eigenvalue''; it does not prove that the repeated
case is hard, nor that geometric simplicity is necessary for polynomial solvability.
